% Where phi functions go in a loop (Lesson 11): the SSA form of a
% counting loop.  B2 is a join of the entry path and the back edge, so
% i and s receive phi functions there; n has one definition and needs
% none; B4 has a single predecessor, so no phi is placed there although
% two definitions of s reach it.
% Original figure (applies the criterion of the outline figure
% m11l2_path_converge to the book's own CFG).
% Usage: % Where phi functions go in a loop (Lesson 11): the SSA form of a
% counting loop.  B2 is a join of the entry path and the back edge, so
% i and s receive phi functions there; n has one definition and needs
% none; B4 has a single predecessor, so no phi is placed there although
% two definitions of s reach it.
% Original figure (applies the criterion of the outline figure
% m11l2_path_converge to the book's own CFG).
% Usage: % Where phi functions go in a loop (Lesson 11): the SSA form of a
% counting loop.  B2 is a join of the entry path and the back edge, so
% i and s receive phi functions there; n has one definition and needs
% none; B4 has a single predecessor, so no phi is placed there although
% two definitions of s reach it.
% Original figure (applies the criterion of the outline figure
% m11l2_path_converge to the book's own CFG).
% Usage: % Where phi functions go in a loop (Lesson 11): the SSA form of a
% counting loop.  B2 is a join of the entry path and the back edge, so
% i and s receive phi functions there; n has one definition and needs
% none; B4 has a single predecessor, so no phi is placed there although
% two definitions of s reach it.
% Original figure (applies the criterion of the outline figure
% m11l2_path_converge to the book's own CFG).
% Usage: \input{figures/l11-path-convergence}
% Width: body (120 mm); the drawing is about 112 mm wide.
\begin{tikzpicture}[x=1mm, y=1mm, font=\scriptsize\sffamily,
    >={Stealth[length=1.6mm]},
    bb/.style={draw, fill=ln-box, align=left, inner sep=2pt,
               minimum width=30mm, rounded corners=0.8mm}]
  \node[bb] (b1) at (22,46) {\textbf{B1}\\\texttt{i1 = 0}\\\texttt{s1 = 0}};
  \node[bb, fill=ln-tip!15] (b2) at (22,27)
    {\textbf{B2}\\\texttt{i2 = $\phi$(i1, i3)}\\\texttt{s2 = $\phi$(s1, s3)}\\\texttt{if (i2 < n)}};
  \node[bb] (b3) at (66,27) {\textbf{B3}\\\texttt{s3 = s2 + i2}\\\texttt{i3 = i2 + 1}};
  \node[bb] (b4) at (22,6) {\textbf{B4}\\\texttt{return s2}};
  \draw[->] (b1) -- (b2);
  \draw[->] (b2.east) -- node[above] {\texttt{i2 < n}} (b3.west);
  \draw[->, ln-caution, thick] (b3.south) -- ++(0,-6) -| 
    node[pos=0.25, below] {\iflangja{後退辺}{back edge}} ([xshift=4mm]b2.south);
  \draw[->] ([xshift=-6mm]b2.south) -- node[left] {\texttt{i2 >= n}} ([xshift=-6mm]b4.north);
  \node[anchor=west, align=left, text=ln-muted, text width=34mm] at (84,40)
    {\iflangja{B2: 2経路（B1 から，B3 から）が B2 でのみ交わる $\Rightarrow$ $\phi$}{B2: the paths from B1 and from B3 meet only at B2 $\Rightarrow$ $\phi$}};
  \node[anchor=west, align=left, text=ln-muted, text width=34mm] at (84,8)
    {\iflangja{B4: 経路はすべて B2 を通る $\Rightarrow$ $\phi$ 不要}{B4: every path passes through B2 $\Rightarrow$ no $\phi$}};
\end{tikzpicture}

% Width: body (120 mm); the drawing is about 112 mm wide.
\begin{tikzpicture}[x=1mm, y=1mm, font=\scriptsize\sffamily,
    >={Stealth[length=1.6mm]},
    bb/.style={draw, fill=ln-box, align=left, inner sep=2pt,
               minimum width=30mm, rounded corners=0.8mm}]
  \node[bb] (b1) at (22,46) {\textbf{B1}\\\texttt{i1 = 0}\\\texttt{s1 = 0}};
  \node[bb, fill=ln-tip!15] (b2) at (22,27)
    {\textbf{B2}\\\texttt{i2 = $\phi$(i1, i3)}\\\texttt{s2 = $\phi$(s1, s3)}\\\texttt{if (i2 < n)}};
  \node[bb] (b3) at (66,27) {\textbf{B3}\\\texttt{s3 = s2 + i2}\\\texttt{i3 = i2 + 1}};
  \node[bb] (b4) at (22,6) {\textbf{B4}\\\texttt{return s2}};
  \draw[->] (b1) -- (b2);
  \draw[->] (b2.east) -- node[above] {\texttt{i2 < n}} (b3.west);
  \draw[->, ln-caution, thick] (b3.south) -- ++(0,-6) -| 
    node[pos=0.25, below] {\iflangja{後退辺}{back edge}} ([xshift=4mm]b2.south);
  \draw[->] ([xshift=-6mm]b2.south) -- node[left] {\texttt{i2 >= n}} ([xshift=-6mm]b4.north);
  \node[anchor=west, align=left, text=ln-muted, text width=34mm] at (84,40)
    {\iflangja{B2: 2経路（B1 から，B3 から）が B2 でのみ交わる $\Rightarrow$ $\phi$}{B2: the paths from B1 and from B3 meet only at B2 $\Rightarrow$ $\phi$}};
  \node[anchor=west, align=left, text=ln-muted, text width=34mm] at (84,8)
    {\iflangja{B4: 経路はすべて B2 を通る $\Rightarrow$ $\phi$ 不要}{B4: every path passes through B2 $\Rightarrow$ no $\phi$}};
\end{tikzpicture}

% Width: body (120 mm); the drawing is about 112 mm wide.
\begin{tikzpicture}[x=1mm, y=1mm, font=\scriptsize\sffamily,
    >={Stealth[length=1.6mm]},
    bb/.style={draw, fill=ln-box, align=left, inner sep=2pt,
               minimum width=30mm, rounded corners=0.8mm}]
  \node[bb] (b1) at (22,46) {\textbf{B1}\\\texttt{i1 = 0}\\\texttt{s1 = 0}};
  \node[bb, fill=ln-tip!15] (b2) at (22,27)
    {\textbf{B2}\\\texttt{i2 = $\phi$(i1, i3)}\\\texttt{s2 = $\phi$(s1, s3)}\\\texttt{if (i2 < n)}};
  \node[bb] (b3) at (66,27) {\textbf{B3}\\\texttt{s3 = s2 + i2}\\\texttt{i3 = i2 + 1}};
  \node[bb] (b4) at (22,6) {\textbf{B4}\\\texttt{return s2}};
  \draw[->] (b1) -- (b2);
  \draw[->] (b2.east) -- node[above] {\texttt{i2 < n}} (b3.west);
  \draw[->, ln-caution, thick] (b3.south) -- ++(0,-6) -| 
    node[pos=0.25, below] {\iflangja{後退辺}{back edge}} ([xshift=4mm]b2.south);
  \draw[->] ([xshift=-6mm]b2.south) -- node[left] {\texttt{i2 >= n}} ([xshift=-6mm]b4.north);
  \node[anchor=west, align=left, text=ln-muted, text width=34mm] at (84,40)
    {\iflangja{B2: 2経路（B1 から，B3 から）が B2 でのみ交わる $\Rightarrow$ $\phi$}{B2: the paths from B1 and from B3 meet only at B2 $\Rightarrow$ $\phi$}};
  \node[anchor=west, align=left, text=ln-muted, text width=34mm] at (84,8)
    {\iflangja{B4: 経路はすべて B2 を通る $\Rightarrow$ $\phi$ 不要}{B4: every path passes through B2 $\Rightarrow$ no $\phi$}};
\end{tikzpicture}

% Width: body (120 mm); the drawing is about 112 mm wide.
\begin{tikzpicture}[x=1mm, y=1mm, font=\scriptsize\sffamily,
    >={Stealth[length=1.6mm]},
    bb/.style={draw, fill=ln-box, align=left, inner sep=2pt,
               minimum width=30mm, rounded corners=0.8mm}]
  \node[bb] (b1) at (22,46) {\textbf{B1}\\\texttt{i1 = 0}\\\texttt{s1 = 0}};
  \node[bb, fill=ln-tip!15] (b2) at (22,27)
    {\textbf{B2}\\\texttt{i2 = $\phi$(i1, i3)}\\\texttt{s2 = $\phi$(s1, s3)}\\\texttt{if (i2 < n)}};
  \node[bb] (b3) at (66,27) {\textbf{B3}\\\texttt{s3 = s2 + i2}\\\texttt{i3 = i2 + 1}};
  \node[bb] (b4) at (22,6) {\textbf{B4}\\\texttt{return s2}};
  \draw[->] (b1) -- (b2);
  \draw[->] (b2.east) -- node[above] {\texttt{i2 < n}} (b3.west);
  \draw[->, ln-caution, thick] (b3.south) -- ++(0,-6) -| 
    node[pos=0.25, below] {\iflangja{後退辺}{back edge}} ([xshift=4mm]b2.south);
  \draw[->] ([xshift=-6mm]b2.south) -- node[left] {\texttt{i2 >= n}} ([xshift=-6mm]b4.north);
  \node[anchor=west, align=left, text=ln-muted, text width=34mm] at (84,40)
    {\iflangja{B2: 2経路（B1 から，B3 から）が B2 でのみ交わる $\Rightarrow$ $\phi$}{B2: the paths from B1 and from B3 meet only at B2 $\Rightarrow$ $\phi$}};
  \node[anchor=west, align=left, text=ln-muted, text width=34mm] at (84,8)
    {\iflangja{B4: 経路はすべて B2 を通る $\Rightarrow$ $\phi$ 不要}{B4: every path passes through B2 $\Rightarrow$ no $\phi$}};
\end{tikzpicture}
